\subsubsection{两条技术路线如何在赛场结合}

视觉组常用两条互补路线：

\begin{itemize}
    \item \textbf{传统图像处理}：阈值、形态学、轮廓、几何约束等，延迟低、可解释性强，适合灯条、颜色分割等结构稳定目标。
    \item \textbf{机器学习 / 深度学习}：检测、分类、分割网络，上限高、适应复杂背景，但需数据、算力与部署（TensorRT 等）。
\end{itemize}

赛场上几乎从不是「二选一」：深度网络出框或类别后，仍常接几何解算、滤波与控制；传统方法也可为网络提供感兴趣区域（Region of Interest，ROI）或校验。
细节分别见卷二 OpenCV 与深度学习章。

\begin{example}[把网络与几何接起来]
在纸上写下“图像 → 网络框 → 框内灯条角点 → 位姿”。逐项标注数据：图像为像素阵列，框和角点为像素坐标，位姿为带坐标系的旋转与平移。检查结果：网络给出候选区域，几何模块还需标定与目标尺寸，检测框本身不能直接当距离。
\end{example}

\begin{exercise}
\label{v3-ex-overview-1}
网络已给出装甲检测框，为什么仍可能需要传统图像处理？写出一种组合方式。
\end{exercise}

\subsubsection{三个方向}

组内方向固定为：

\begin{itemize}
    \item \textbf{自瞄}：车载相机识别装甲，估计与预测目标，驱动云台达到可击打姿态。
    \item \textbf{雷达}：高台俯视感知全场，融合（可选）无线电等信息，经裁判系统向队友广播。
    \item \textbf{导航}：移动平台（主要是哨兵）定位、建图与路径规划，自主行进与任务执行。
\end{itemize}

「识别」「CUDA」「激光雷达传感器」等是\textbf{技能或器件}，不是与上述三者并列的方向名。
工具链与部署支撑三个方向，但不单独算作第四方向。

\begin{example}[给需求分配方向]
将“跟随装甲转云台”“绕开障碍到补给点”“共享全场目标位置”写成三张卡，分别放到自瞄、导航、雷达栏。再把 GPU 加速贴在三个栏下方，表示可共同使用的技术。检查：需求按输出对象分类，器件不单独成为方向。
\end{example}

\begin{exercise}
\label{v3-ex-overview-2}
把“用激光雷达定位”“用网络识别装甲”“把敌方位置广播给队友”分别归类；哪项仅凭描述不能唯一归类？
\end{exercise}

\subsubsection{视觉数据流（总览）}

典型路径可概括为：

\begin{figure}[H]
\centering
\resizebox{0.9\textwidth}{!}{%
\begin{tikzpicture}[box/.style={draw=cyan!60!black,rounded corners,align=center,
  minimum height=10mm,text width=27mm,fill=cyan!5},>=stealth]
\node[box] (sensor) at (0,0) {相机 / 传感器\\图像、量测};
\node[box] (state) at (4.2,0) {识别、估计\\定位或融合};
\node[box] (aim) at (8.4,1.6) {自瞄\\目标角与有效性};
\node[box] (nav) at (8.4,0) {导航\\路径与速度};
\node[box] (radar) at (8.4,-1.6) {雷达\\场地坐标广播};
\draw[->,thick] (sensor) -- (state);
\draw[->,thick] (state.east) -- (aim.west);
\draw[->,thick] (state) -- (nav);
\draw[->,thick] (state.east) -- (radar.west);
\end{tikzpicture}}
\caption{按输出职责区分三个方向；中间处理模块按任务选择}
\label{fig:v3-dataflow}
\end{figure}

\begin{itemize}
    \item 自瞄走到「云台控制」闭环；
    \item 导航走到「速度 / 路径」；
    \item 雷达走到「广播与辅助决策」，一般不直接拧本车云台。
\end{itemize}

\begin{example}[追踪一条观测]
在纸上按顺序填写：相机产生第 42 帧，检测输出像素角点，解算输出相机系位置，跟踪输出带时间戳的状态，控制接口输出目标角。给每个箭头写上数据名称。检查：图像坐标不能直接送入要求角度的接口。
\end{example}

\begin{exercise}
\label{v3-ex-overview-3}
雷达检测得到像素点后，能否直接把这两个数当作地图坐标广播？指出缺少的步骤。
\end{exercise}

\subsubsection{与其它组别的接口边界}

\begin{itemize}
    \item \textbf{机械}：相机位姿、云台限位、线束与减振；视觉提出视场（Field of View，FOV）与刚度需求，不代替结构设计。
    \item \textbf{电控}：云台/底盘控制接口、裁判系统串口、时间同步约定；视觉给目标角或速度建议，执行器闭环多在电控。
    \item \textbf{运营 / 战术}：雷达与导航的信息如何服务于决策；视觉保证数据时效与协议正确。
\end{itemize}

接口变更必须书面同步（消息字段、单位、坐标系、频率），禁止「口头改了但仓库未改」。

\begin{example}[一张最小接口卡]
以教学接口为例，写下：发送方视觉、接收方电控；yaw 为相对当前云台的偏航增量，单位 rad，正方向由双方坐标约定决定；另含采样时刻与有效标志。让接收方按卡片解释正负号与失效处理。双方解释一致才算填写完成；具体字段名以队内协议为准。
\end{example}

\begin{note}
接口必须说明单位、坐标系、时间语义和失效行为。同名字段不代表双方理解一致，变更后用已知输入验证。
\end{note}

\begin{conclusion}
\begin{itemize}
  \item 按任务输出区分自瞄、雷达、导航。
  \item 像素、空间位置、控制量各有单位和坐标约定。
  \item 工具与器件可被多个方向复用。
\end{itemize}
\end{conclusion}

\begin{exercise}
\label{v3-ex-overview-4}
同一个 yaw 字段从度改成弧度，至少要同步哪些信息？
\end{exercise}

\subsubsection{本卷读法}

\begin{center}
\begin{tabular}{@{}p{0.30\textwidth}p{0.62\textwidth}@{}}
\toprule[1.5pt]
问题 & 去哪 \\
\midrule[1pt]
不会装环境 / CMake / C++ & 卷一 \\
不会写节点、不懂 PnP/Kalman 公式 & 卷二 \\
不知道链路怎么串、调车守什么规矩 & \textbf{本卷} \\
\bottomrule[1.5pt]
\end{tabular}
\end{center}
开营建议先读第\ref{ch:v3-overview}、\ref{ch:v3-aim}、\ref{ch:v3-engineering}章；相机上机前读第\ref{ch:v3-camera}章。通识过关后按方向选修卷二；中期用第\ref{ch:v3-aim}--\ref{ch:v3-radar}章做链路答辩。

\begin{example}[查找下一本书]
以“会识别却不懂为什么位置跳变”为例：先用本卷标出检测、解算、估计三段，再到卷二查相机标定、几何视觉与滤波。记录章名及要找的输入输出，而非只写“去看算法”。
\end{example}

\begin{exercise}
\label{v3-ex-overview-5}
编译报错、PnP 输入点顺序错误、接口坐标系未约定，分别优先去哪里查？
\end{exercise}
